
\medskip
\phantomsection\label{thm:sga3-VIB-6.5.2}%
\noindent\textbf{Theorem 6.5.2.---}
\emph{Let \(X\) be a smooth scheme over \(S\), with geometrically
irreducible fibers.}%
\begin{enumerate}
\item[\textup{(i)}]\phantomsection\label{concl:sga3-VIB-6.5.2-i}
  \emph{For every closed subscheme \(Y\) of \(X\), the functor
  \[
    \prod_{X/S}Y/X
  \]
  is representable by a closed subscheme \(T\) of \(S\).}
\item[\textup{(ii)}]\phantomsection\label{concl:sga3-VIB-6.5.2-ii}
  \emph{Moreover, if \(Y\to X\) is of finite presentation, then
  so is \(T\to S\).}
\end{enumerate}

\emph{Proof.}
Let \(f:X\to S\) denote the structure morphism. Since \(f\) is
faithfully flat and locally of finite
presentation, it is a covering for the fpqc topology. Moreover,
\[
  T=\prod_{X/S}Y/X
\]
is evidently a subsheaf of \(S\) for the fpqc topology. It follows
that the representability of \(T\) by a closed subscheme of \(S\)
is a question local on \(S\) for the fpqc topology,%
and the same is true of the question whether \(T\) is of finite
presentation over \(S\) (EGA~IV\(_2\), 2.7.1).

After the base change \(S'\to S\) with \(S'=X\), we are reduced
to the case in which \(X\) admits a section
\[
  \epsilon:S\longrightarrow X.
\]
We may also suppose that \(S\) is affine, and hence
quasi-compact. We then have:

\medskip
\phantomsection\label{cor:sga3-VIB-6.5.3}%
\noindent\textbf{Corollary 6.5.3.---}
\emph{Under the hypotheses of
\hyperref[thm:sga3-VIB-6.5.2]{6.5.2}, suppose that \(S\) is
quasi-compact, that \(X\to S\) admits a section \(\epsilon\), and
that \(Y\to X\) is of finite presentation. Then there exists an
integer \(n\geq0\) such that}
\phantomsection\label{eq:sga3-VIB-6.5.3-infinitesimal}%
\[
  \prod_{X/S}Y/X
  =
  \prod_{X_n/S}Y_n/X_n,
\]
\emph{where \(X_n\) is the \(n\)-th infinitesimal neighborhood of
the immersion \(\epsilon:S\to X\), and
\(Y_n=Y\cap X_n\).}

When \(Y\) is of finite presentation over \(X\), this corollary
does indeed imply
\hyperref[thm:sga3-VIB-6.5.2]{6.5.2}. Since \(X\) is smooth over
\(S\), \(X_n\) is finite and locally free over \(S\), and is
therefore, a fortiori, essentially free over \(S\)
(see \hyperref[def:sga3-VIB-6.2.1]{6.2.1}). Hence
\[
  \prod_{X_n/S}Y_n/X_n
\]
is representable by a closed subscheme \(T\) of \(S\), of finite
presentation over \(S\), by
\hyperref[thm:sga3-VIB-6.2.3]{6.2.3}.

Let us first prove \hyperref[cor:sga3-VIB-6.5.3]{6.5.3}, and
hence \hyperref[thm:sga3-VIB-6.5.2]{6.5.2}, when \(S\) is
Noetherian. Put
\phantomsection\label{eq:sga3-VIB-6.5.3-Tn}%
\[
  T_n=\prod_{X_n/S}Y_n/X_n.
\]
The \(T_n\) form a decreasing sequence of closed subschemes of
\(S\), and since \(S\) is Noetherian this sequence is stationary.
Let
\phantomsection\label{eq:sga3-VIB-6.5.3-R}%
\[
  R=\bigcap_{n\geq0}\prod_{X_n/S}Y_n/X_n
\]
be their common value for all sufficiently large \(n\). Clearly
\(T\subset R\), and it is enough to prove \(R\subset T\).

After the base change \(R\to S\), we are reduced to the case
\[
  R=S,
\]
that is,
\[
  Y_n=X_n\quad\text{for every \(n\)},
  \qquad\text{or equivalently}\qquad
  Y\supset X_n\quad\text{for every \(n\)}.
\]
We must then prove \(T=S\), that is, \(Y=X\).

Now \(Y\supset X_n\) for every \(n\), together with the fact that
\(X\) is locally Noetherian, implies that in a neighborhood of
each point of \(\epsilon(S)\), \(Y\) is an induced open subscheme
of \(X\).%
Consequently, there exists an induced open subscheme \(U\) of
\(X\), containing \(\epsilon(S)\), such that \(U\subset Y\). By
EGA~IV\(_3\), 11.10.10, since the fibers of \(X/S\) are integral,
\(U\) is schematically dense in \(X\). Since \(Y\) is a closed
subscheme containing \(U\), it follows that \(Y=X\). This proves
\hyperref[cor:sga3-VIB-6.5.3]{6.5.3}, and hence
\hyperref[thm:sga3-VIB-6.5.2]{6.5.2}, when \(S\) is Noetherian.

The general case is reduced to the preceding one. For every
\(s\in S\), there exist an affine open neighborhood \(U\) of
\(s\) and an affine open neighborhood \(V\) of \(\epsilon(s)\)
such that \(f(V)\subset U\). Then \(f(V)\) is an open
neighborhood of \(s\) contained in \(U\). Let \(S'\) be an affine
open neighborhood of \(s\) contained in
\[
  \epsilon^{-1}(V)\cap f(V),
\]
and put
\[
  X'=V\cap f^{-1}(S').
\]
Then \(X'\) and \(S'\) are affine open subsets of \(X\) and \(S\),
respectively, and \(X'/S'\) admits a section. By the local nature
of \SGARefLink{thm:sga3-VIB-6.5.2}{6.5.2} and \SGARefLink{cor:sga3-VIB-6.5.3}{6.5.3}, we may suppose that \(S=S'\).

Thus \(X'\) is an affine open subset of \(X\) containing
\(\epsilon(S)\). Since every fiber \(X_s\) is irreducible,
\(X'_s\) is schematically dense in \(X_s\). Consequently, by
EGA~IV\(_3\), 11.10.10, \(X'\) is schematically dense in \(X\);
likewise, after every base change \(S_1\to S\),
\[
  X'\times_S S_1
\]
is schematically dense in \(X\times_S S_1\).

It follows that
\phantomsection\label{eq:sga3-VIB-6.5.3-open-restriction}%
\[
  \prod_{X/S}Y/X
  =
  \prod_{X'/S}Y'/X',
  \qquad Y'=Y\cap X'.
\]
This reduces us to the case \(X=X'\), so we may suppose that both
\(S\) and \(X\) are affine. Moreover, if
\[
  X=\operatorname{Spec}B
\]
and \(J\) is the ideal of \(B\) defining \(Y\), then \(J\) is the
filtered colimit of its finitely generated subideals. Thus \(Y\)
is the intersection of closed subschemes \(Y_i\) of \(X\) that
are of finite presentation over \(X\), and consequently
\phantomsection\label{eq:sga3-VIB-6.5.3-finite-presentation-reduction}%
\[
  \prod_{X/S}Y/X
  =
  \bigcap_i\prod_{X/S}Y_i/X.
\]
To prove \hyperref[thm:sga3-VIB-6.5.2]{6.5.2}, we are therefore
reduced to the case in which \(Y\) is of finite presentation over
\(X\), with \(S\) and \(X\) affine.

In that case \(X\) and \(Y\) over \(S\) are obtained by base
change \(S\to S_0\) from an analogous situation \(X_0,Y_0\) over
a Noetherian scheme \(S_0\). This reduces us to the Noetherian
case already treated, and completes the proof of
\hyperref[thm:sga3-VIB-6.5.2]{6.5.2} and
\hyperref[cor:sga3-VIB-6.5.3]{6.5.3}.
